% Moved from sections/13_algebra.tex (end-sem split: questions live in the Question Bank PDF).
\section{Algebraic Structures: Groups, Rings, Fields \hothot}\label{q:13_algebra}
\noindent{\small Notes: \S\ref{sec:algebra}.}

\subsection*{Check your understanding}
\begin{enumerate}[leftmargin=*]
  \item \textbf{Is $\langle \mathbb{Z}_n^*, \times \rangle$ a group? Is $\langle \mathbb{Z}_n, \times \rangle$?} $\mathbb{Z}_n^*$ is an abelian group: it is closed (a product of numbers coprime to $n$ is coprime to $n$), associative, has identity 1, and every element has an inverse by extended Euclid because $\gcd(a, n) = 1$. $\mathbb{Z}_n$ under $\times$ is not a group, because 0 (and every $a$ with $\gcd(a, n) > 1$) has no inverse.
  \item \textbf{Is $\mathbb{Z}$ a field?} No. $\langle \mathbb{Z}, +, \times \rangle$ is a commutative ring: addition satisfies all five properties, but multiplication has no inverses except for $\pm 1$, so division leaves the set ($12 \div 5 = 2.4 \notin \mathbb{Z}$).
  \item \textbf{Why is $\mathbb{Z}_{256}$ not a field, and what does AES use instead?} 256 is not prime, so there are zero divisors: $2 \times 128 \equiv 0$, which means 2 (and every even byte) has no inverse. AES uses GF$(2^8)$, where bytes are polynomials, addition is XOR and multiplication is reduced modulo the irreducible $x^8 + x^4 + x^3 + x + 1$, so all 256 bytes are field elements.
  \item \textbf{Differentiate a group, a ring and a field.} A group has one operation with closure, associativity, identity and inverses, and so supports one pair such as $(+, -)$; a ring adds a second operation that only needs closure and associativity and distributes over the first, so it supports $(+, -, \times)$; a field also requires inverses for the second operation (except for 0), so it supports $+, -, \times, \div$. Examples: $\langle \mathbb{Z}_n, + \rangle$, $\langle \mathbb{Z}, +, \times \rangle$, $\langle \mathbb{Z}_p, +, \times \rangle$.
  \item \textbf{Why does applying two transposition ciphers one after another not make the cipher stronger?} Permutations of the same block length form a group under composition, so by closure the two permutations equal one permutation $\pi_2 \circ \pi_1$, and the attacker only has to find that one key from the same key space of $k!$.
  \item \textbf{Is $\langle \mathbb{Z}_{10}, + \rangle$ a subgroup of $\langle \mathbb{Z}_{12}, + \rangle$?} No. Although $\mathbb{Z}_{10} \subset \mathbb{Z}_{12}$ as sets, a subgroup must use the same operation, and addition mod 10 is not addition mod 12.
  \item \textbf{Can $\mathbb{Z}_{11}^*$ contain an element of order 4?} No. $|\mathbb{Z}_{11}^*| = 10$, and by Lagrange the order of an element divides the order of the group, so only the orders 1, 2, 5 and 10 are possible.
  \item \textbf{Why must the modulus of GF$(2^n)$ be irreducible? Can $x^4 + x^2 + 1$ be used for GF$(2^4)$?} If $m = f \cdot g$, then $f \otimes g \equiv 0$ with $f, g \ne 0$, so $f$ has no inverse and the structure is not a field. $x^4 + x^2 + 1 = (x^2 + x + 1)^2$ is reducible even though it has no root, so it cannot be used.
  \item \textbf{Why are addition and subtraction the same in GF$(2^n)$?} The coefficients are in GF(2), where $-1 \equiv 1$, so $-B = B$ and $A - B = A + B$; bit by bit, both are XOR.
  \item \textbf{Find $\{02\} \bullet \{C3\}$ and $x^{-1}$ in GF$(2^8)$.} \ttx{C3} $=$ \ttx{11000011} has its top bit set, so shift to \ttx{10000110} and XOR \ttx{1B}: \textbf{\ttx{9D}}. For $x^{-1}$: $x^8 \equiv x^4 + x^3 + x + 1$ means $x \cdot (x^7 + x^3 + x^2 + 1) \equiv 1$, so $x^{-1} = x^7 + x^3 + x^2 + 1$ (\ttx{8D}).
  \item \textbf{In GF(11), find $3 \div 7$.} $7^{-1} = 8$ since $56 = 5 \cdot 11 + 1$, so $3 \div 7 = 3 \times 8 = 24 \equiv \mathbf{2}$; check $2 \times 7 = 14 \equiv 3$.
\end{enumerate}
